\section{Campana--Peternell and general Schnell fiber spaces}
\label{sec:general-consequences}

The good-model input supplies canonical nonvanishing as well as the
model used in the mixed-intersection argument. These two outputs fit
the reduction of \cite[Sections~4--10]{Schnell2022}.

\subsection{Nonvanishing and very general fibers}

Let $T$ be a smooth connected projective complex variety with
pseudo-effective $K_T$, and take the model $V$ and comparison divisor
$E$ supplied by Theorem~\ref{thm:canonical-good-models}.
Choose a positive multiple $rK_V$ that is Cartier and globally
generated. The comparison then makes $rE$ an effective integral
Cartier divisor. The section transfer
in~\eqref{eq:section-transfer} gives
\[
 0\ne H^0(V,\cO_V(rK_V))
 \lhook\joinrel\longrightarrow H^0(T,\cO_T(rK_T)).
\]
Consequently,
\begin{equation}\label{eq:canonical-nonvanishing}
 K_T\in\PE(T)\quad\Longrightarrow\quad \kappa(T)\geq0.
\end{equation}
This is the canonical nonvanishing input used by Schnell.

We will use smooth closed complex fibers when applying
\eqref{eq:canonical-nonvanishing}. They have the same Kodaira
dimension as the geometric generic fiber when chosen very generally.
Indeed, let $f\colon X\to Y$ be a surjective morphism with connected
fibers between smooth connected projective complex varieties, and
write $X_{\bar\eta}$ for its geometric generic fiber. Over a nonempty
open subset, $f$ is smooth; there
$\omega_{X/Y}=\omega_X\otimes f^*\omega_Y^{-1}$ restricts to the
canonical bundle of each fiber. For each positive integer $m$, after
shrinking this open subset, formation of
$f_*(\omega_{X/Y}^{\otimes m})$ commutes with base change.
Outside the union of the resulting countably many proper
closed subsets, a smooth fiber $F$ therefore satisfies
\[
 h^0(F,\omega_F^{\otimes m})
 =
 h^0(X_{\bar\eta},\omega_{X_{\bar\eta}}^{\otimes m})
 \qquad\text{for every }m>0.
\]
Here the dimension on the right is over the algebraic closure of
$\C(Y)$. Thus the plurigenus sequences, and hence the Kodaira
dimensions, agree. Such complex fibers exist because $\C$ is
uncountable. In particular, the zero-Kodaira fiber condition in the
reduction below is precisely the geometric generic fiber condition
of Theorem~\ref{thm:schnell}.

We also use the restriction observation in
\cite[Section~7]{Schnell2022}. A pseudo-effective real Cartier
divisor $L$ on $X$ restricts to a pseudo-effective divisor on a very
general fiber. To see this, choose effective real divisors whose
numerical classes converge to $[L]$. After excluding a countable
union of proper closed subsets of $Y$, a fiber is contained in none
of their supports. Their restrictions are effective, and their
numerical classes converge to the class of $L|_F$. Applying this to
$L=m_0K_X-f^*H$ gives
\[
 K_F\in\PE(F),\qquad \kappa(F)\geq0,
\]
where the second assertion follows from
\eqref{eq:canonical-nonvanishing}. The very general fiber can be
chosen to satisfy this restriction property and the plurigenus
comparison simultaneously.

\subsection{The Campana--Peternell inequality}

\begin{proof}[Proof of Corollary~\ref{cor:campana-peternell}]
Since $D$ is effective, $\kappa(D)\geq0$. Also
\[
 K_X=\frac1{m_0}\bigl((m_0K_X-D)+D\bigr)\in\PE(X).
\]
If $\kappa(D)=0$, canonical nonvanishing
\eqref{eq:canonical-nonvanishing} proves the assertion. We may
therefore assume $\kappa(D)>0$.

Schnell's reduction in \cite[Sections~4--6]{Schnell2022} first
resolves a system $|nD|$ whose image has dimension $\kappa(D)$,
then takes its Stein factorization and resolves the base. It produces
an algebraic fiber space
$f_0\colon X_0\to Y_0$ between smooth connected projective complex
varieties, with $X_0$ birational to $X$, an ample Cartier divisor
$H_0$ on $Y_0$, and a positive integer $a_0$ such that
\begin{equation}\label{eq:initial-cp-reduction}
 \dim Y_0=\kappa(D),\qquad
 a_0K_{X_0}-f_0^*H_0\in\PE(X_0).
\end{equation}
The effective exceptional terms from resolving $X$ preserve
pseudo-effectivity; after resolving the base, Schnell replaces the
big and nef pullback of the original ample divisor by a suitable
ample divisor. In particular, the output has the precise ample
Cartier hypothesis of Theorem~\ref{thm:schnell}.

We recall the finite iteration in
\cite[Sections~7 and~9]{Schnell2022}. Suppose that
$f_i\colon X_i\to Y_i$ is such a fiber space, with ample Cartier
$H_i$ and
$a_iK_{X_i}-f_i^*H_i\in\PE(X_i)$ for a positive integer $a_i$.
For a very general smooth fiber $F_i$, the preceding restriction and
nonvanishing argument gives $\kappa(F_i)\geq0$. If
$\kappa(F_i)>0$, the construction in
\cite[Lemma~7.1]{Schnell2022} chooses positive integers $r_i,b_i$
and an effective Cartier divisor $L_i$ such that
\[
 L_i\sim r_iK_{X_i}+f_i^*(b_iH_i),\qquad
 \kappa(L_i)=\kappa(F_i)+\dim Y_i.
\]
The resulting comparison is
\[
 (a_ib_i+r_i)K_{X_i}-L_i
 \sim b_i(a_iK_{X_i}-f_i^*H_i).
\]
Its right side is pseudo-effective, so the same is true of its left
side, since linear equivalence preserves numerical classes.
Applying the divisor-to-fiber-space reduction of Sections~4--6 of
Schnell's paper to $L_i$ gives another algebraic fiber space
$f_{i+1}\colon X_{i+1}\to Y_{i+1}$ of the same smooth projective
complex type, with $X_{i+1}$ birational to $X_i$, an ample Cartier
divisor $H_{i+1}$, and a positive integer $a_{i+1}$ such that
\[
 \dim Y_{i+1}=\kappa(F_i)+\dim Y_i,\qquad
 a_{i+1}K_{X_{i+1}}-f_{i+1}^*H_{i+1}\in\PE(X_{i+1}).
\]

As long as $\kappa(F_i)>0$, the integer $\dim Y_i$ strictly
increases. It is bounded by $\dim X$, so this process reaches an
index $j$ with $\kappa(F_j)=0$. The geometric generic fiber has
Kodaira dimension zero by the comparison above.
Theorem~\ref{thm:schnell} applies at this last stage. Birational
invariance of Kodaira dimension and
\eqref{eq:initial-cp-reduction} give
\[
 \kappa(X)=\kappa(X_j)=\dim Y_j
 \geq\dim Y_0=\kappa(D).
\]
\end{proof}

If $D$ is an effective rational divisor instead, choose a positive
integer $q$ such that $qD$ is an integral Cartier divisor.
Multiplying the numerical hypothesis by $q$ gives
$qm_0K_X-qD\in\PE(X)$. Corollary~\ref{cor:campana-peternell}
applied to $qD$ gives the same conclusion because
$\kappa(qD)=\kappa(D)$. This extension still requires an effective
divisor; it makes no assertion for an arbitrary pseudo-effective
divisor in its place.

\subsection{Equality and eventual effectivity for general fibers}

\begin{proof}[Proof of Corollary~\ref{cor:schnell-general}]
If $Y$ is a point, then $F=X$ and $H$ is linearly equivalent to zero.
The equality is tautological, and the numerical hypothesis gives
pseudo-effective $K_X$. A nonzero section supplied by
\eqref{eq:canonical-nonvanishing}, together with its powers, gives
the asserted sections.

Assume $\dim Y>0$. The restriction argument above gives
$\kappa(F)\geq0$. As in \cite[Lemma~7.1]{Schnell2022}, choose
positive integers $a,b$ and an effective Cartier divisor $L$ with
\[
 L\sim aK_X+f^*(bH),\qquad \kappa(L)=\kappa(F)+\dim Y.
\]
The linear equivalence
\[
 (m_0b+a)K_X-L\sim b(m_0K_X-f^*H)
\]
shows that its left side is pseudo-effective.
Corollary~\ref{cor:campana-peternell}, applied to the effective
Cartier divisor $L$, gives
\[
 \kappa(X)\geq\kappa(L)=\kappa(F)+\dim Y.
\]
The reverse inequality is the easy-addition inequality used in
\cite[Lemma~7.1]{Schnell2022}. This proves the equality.

To obtain sections after subtracting the original $f^*H$, we use
the Fujita--Mori criterion recalled in
\cite[Lemma~4.6]{PopaSchnell2014}; this is the criterion behind
\cite[Section~8]{Schnell2022}. Since $\kappa(F)\geq0$, the equality
just proved implies that there are a big Cartier divisor $B$ on $Y$
and a positive integer $u$ with a nonzero section
\[
 \sigma\in H^0\bigl(X,\cO_X(uK_X-f^*B)\bigr).
\]
Bigness of $B$ gives a positive integer $v$ and a nonzero section
$\tau\in H^0(Y,\cO_Y(vB-H))$. Thus, with $r=uv$,
\[
 s=\sigma^v f^*\tau
 \in H^0\bigl(X,\cO_X(rK_X-f^*H)\bigr)
\]
is nonzero. This step replaces the big divisor supplied by the
criterion with the ample divisor in the original hypothesis.

Finally, ampleness of $H$ supplies a nonzero section
$\tau_\ell\in H^0(Y,\cO_Y((\ell-1)H))$ for every sufficiently
large integer $\ell$. The products
\[
 s^\ell f^*\tau_\ell
 \in H^0\bigl(X,\cO_X(\ell rK_X-f^*H)\bigr)
\]
are nonzero. Choosing $\ell_0$ beyond this ampleness threshold proves
the asserted conclusion for every integer $\ell\geq\ell_0$.
\end{proof}
